\section{BIRD Admission-Policy Study}\label{app:phase2-protocol}

The study uses a $2{\times}2$ factorial design plus an open-mechanism
arm, with six builders and three seeds. Each slot allows at most
$R=3$ attempts, stopping at first admission without top-up. Two development
databases are excluded from the evaluation pool. The evaluation set
contains $1{,}169$ questions from nine databases and a pre-drawn
$400$-item core. An $18$-item pipeline smoke test is a disclosed
exception to development isolation. The official BIRD judge enforces
a $30$\,s timeout. Responses are cached by model identity with
cross-process locking; a cache-concurrency defect found during
collection was corrected.

The original primary contrast was D$-$A on paired builder--seed cells
with hierarchical bootstrap and permutation inference, with
Holm-corrected factorial secondary analyses on the common core and
K-matched sensitivity analyses. Instrument-calibration fixes preceded
confirmatory outcomes. Other deviations included operational access
handling and quarantine of inconsistent observations. The corrected
analysis uses shared baseline records and database/task resampling;
these revisions followed the original analysis specification.
The implemented gate structurally rejects two permitted prompt-level
strategies, and realized generation costs differ across arms.

\begin{figure}[H]
\centering
\includegraphics[width=0.75\linewidth]{figures/fig_design_revision.pdf}
\caption{\textbf{Phase-II confirmatory design} (frozen protocol). Left:
database-level contamination isolation. Middle: the $2{\times}2$
factorial with a common attempt cap, plus open-mechanism arm II-E. Right:
the three-layer admission pipeline.}
\label{fig:design}
\end{figure}

\subsection{Design and primary contrasts}
\label{sec:f2-admission}
A separate confirmatory BIRD study crosses six independent builders with
three seeds under strategy forcing (a named mechanism from a fixed set
versus free invention) and conformance gating, giving four arms:
II-A (free, ungated), II-B (free, gated), II-C (forced, ungated), and
II-D (forced, gated). The evaluation covers
$1{,}169$ questions from nine databases never opened during
development (one disclosed $18$-item smoke-test exception), with the
official BIRD execution scorer and a pre-outcome-frozen analysis later
re-run as a versioned corrected reanalysis. The implemented gate bundles
nontrivial-mechanism structural screening with conformance verification
--- two of eight forced prompt-transformation slots cannot pass even
when neutrally valid --- so the treatment does not isolate verification
causally.

\begin{table}[H]
\centering\small
\begin{tabular}{@{}lrr@{}}
\toprule
Metric & II-A & II-D \\
\midrule
Candidate count & 5.78 & 3.83 \\
Oracle accuracy (\%) & 76.14 & 75.70 \\
Best-fixed accuracy (\%) & 67.58 & 67.30 \\
Headroom (pp) & 8.56 & 8.40 \\
\midrule
D$-$A headroom (pp) & \multicolumn{2}{c}{$\RevisionPrimaryDelta$} \\
Unadjusted bootstrap 95\% CI & \multicolumn{2}{c}{$\RevisionPrimaryCI$} \\
One-sided sign-flip $p$ & \multicolumn{2}{c}{$\RevisionPrimaryP$} \\
\bottomrule
\end{tabular}

\caption{\textbf{Corrected primary analysis.} Means over $18$
builder--seed cells on $1{,}169$ questions; \bare{} is included in both
the oracle and best-fixed comparators.}
\label{tab:primary}
\end{table}

The diagonal forced/gated minus free/ungated contrast in single-call-inclusive
oracle headroom is $\RevisionPrimaryDelta$\,\%, with cell contrasts
ranging over $\RevisionPrimaryRange$\,\%: the bundled screening policy
does not improve mean headroom. The original uncertainty analysis is
retained in Table~\ref{tab:primary}; this mean comparison does not
establish equivalence. Structural admission is therefore not evidence
that a mechanism activates on relevant tasks or adds stable utility.

\begin{figure}[H]
\centering
\includegraphics[width=0.72\linewidth]{figures/fig_results_revision.pdf}
\caption{\textbf{Versioned corrected reanalysis.} Left: full-set
D$-$A cell contrasts with the pooled 95\% interval. Right: four-arm
headroom on the shared core-$400$ tasks with a common \bare{} baseline.}
\label{fig:results}
\end{figure}

On the shared core-$400$ tasks the gate contrast is
$\RevisionGateDelta$\,\%, and the strategy and interaction secondaries
are likewise inconclusive (Table~\ref{tab:factorial}); the fixed-$K$
mean contrast $\RevisionKmatchedDelta$\,\% provides a population-size
sensitivity check. The open-mechanism arm ran only on
the development side and has no confirmatory estimand. These are all
single-run contrasts, and the execution variability reported in
the R2/R3 audits in Appendix~\ref{app:r2} ($8.54\%$ R2 verdict flips; $14.5\%$ of tasks
changing \bare{} correctness across three runs; a $6.50$\,\% same-code
one-run gap) characterizes them as coverage opportunities rather than
capability differences.

\begin{table}[H]
\centering\small
\begin{tabular}{@{}lrrrr@{}}
\toprule
Contrast & Effect (pp) & 95\% CI (pp) & Raw $p$ & Holm $p$ \\
\midrule
Gate & $-0.69$ & $[-1.56, +0.20]$ & 0.026 & 0.079 \\
Strategy & $+0.42$ & $[-0.32, +1.51]$ & 0.263 & 0.428 \\
Interaction & $-0.87$ & $[-2.83, +0.65]$ & 0.214 & 0.428 \\
\bottomrule
\end{tabular}

\caption{\textbf{Corrected core-$400$ factorial reanalysis.} Gate is
$[(B{-}A)+(D{-}C)]/2$, strategy $[(C{-}A)+(D{-}B)]/2$, interaction
$(D{-}C)-(B{-}A)$; raw and Holm-adjusted $p$-values.}
\label{tab:factorial}
\end{table}



\FloatBarrier
\Needspace{0.25\textheight}
\subsection{Same-core headroom decomposition}\label{app:decomp}
Table~\ref{tab:decomp} separates coverage, best-fixed accuracy, and
headroom using the same core tasks and baseline records in every arm.
\begin{table}[H]
\centering\footnotesize
\begin{tabular}{@{}lrrrrrr@{}}
\toprule
Arm & $K_{\rm cand}$ & $K_{\rm total}$ & Single-call (\%) & Best (\%) & Oracle (\%) & $H$ (\%) \\
\midrule
A & 5.78 & 6.78 & 63.25 & 66.00 & 74.51 & 8.51 \\
B & 6.11 & 7.11 & 63.25 & 66.12 & 74.39 & 8.26 \\
C & 5.61 & 6.61 & 63.25 & 65.64 & 75.01 & 9.38 \\
D & 3.83 & 4.83 & 63.25 & 65.65 & 73.90 & 8.25 \\
\bottomrule
\end{tabular}

\caption{\textbf{Same-core headroom decomposition.} All arms use identical
400 tasks and identical single-call outcomes. $K_{\rm cand}$ excludes single-call and
$K_{\rm total}=K_{\rm cand}+1$; counts and accuracies are cell means.
Headroom is oracle minus best-fixed accuracy, not realized routing gain.}
\label{tab:decomp}
\end{table}


\FloatBarrier
\subsection{Sensitivity, repeat audits, and recorded cost}
\paragraph{Corrected K-matched and R2 replays.}\label{app:r2}
Both replays use the same deterministic record ordering and shared
database/within-database task draws as the corrected primary analysis,
with seeds resampled within fixed builders (10{,}000 replicates).
For K-matching, each A/B and C/D pair uses the smaller candidate count;
every uniform candidate subset is enumerated exactly. Single-call stays in every
subset, and best-fixed is recomputed inside each subset on every task draw.
This is a size sensitivity analysis, not a causal mediation decomposition.
Retaining all 18 cells, including two with singleton A subsets, changes
subset D$-$A from \RevisionRtwoCached\ to \RevisionRtwoOff\,\%
(change \RevisionRtwoChange\,\%, 95\% CI \RevisionRtwoChangeCI). R2
subset headroom excludes single-call because R2 did not acquire a matched single-call
run, so it is distinct from the single-call-inclusive primary estimand; the old
16-cell comparison is reconciled here. These are instability measurements,
not standard errors or a detectability floor for the mean treatment effect;
they do not establish that every contrast is unchanged, or identify what
fraction of between-harness disagreement is noise.

R2 acquired 70 harnesses on 400 tasks, selecting the first two callable
admitted candidates per A/D cell where available. The two cache-off
acquisition groups are kept separate despite a duplicated repeat label;
program identity is checked against the cached runs.
No matched single-call was acquired, so R2 headroom excludes single-call in both
conditions. The corrected analysis retains singleton cells (their subset
headroom is zero), yielding all 18 paired cells.
There are \RevisionRtwoFlips/28{,}000 verdict flips against cached outcomes
(\RevisionRtwoFlipPercent\%) and
\RevisionRtwoRepeatFlips/14{,}400 flips between the two D cache-off passes
(\RevisionRtwoRepeatPercent\%).
For reconciliation only, restricting to the 16 cells with at least two
members in both arms gives D$-$A $\RevisionRtwoLegacyCached$\,\% cached
and $\RevisionRtwoLegacyOff$\,\% cache-off. This explains the difference
from the full-cell estimates; neither analysis identifies future execution
variance, the effect of caching alone, or a stable-complementarity advantage
over clones. Acquisition time and caching condition are not isolated.

\paragraph{R3 single-call-only repeat audit.}\label{app:r3}
R3 is a separate, smaller instability audit: three independent cache-off
executions of the \emph{single-call baseline only} (no generated harnesses) on the
same core-400 tasks, same frozen target (GLM-5.3-Flash, temperature 0).
58 of 400 tasks (14.5\%) changed single-call correctness in at least one of the
three repeats; pairwise flip rates between repeat pairs are 39/400, 40/400,
and 37/400 (9.75\%, 10.0\%, 9.25\%). Treating the three single-call repeats as
three pseudo-members (one execution each) gives a one-run oracle of
$72.25\%$ against a best-single-repeat accuracy of $65.75\%$, a
descriptive $6.50$\,\% same-code oracle gap computed from the recorded
official-scorer outcomes.
R3 therefore measures the observed
instability rate of the baseline itself under temperature-0 API
nondeterminism. This is not a standard error or an effect-size floor for
any contrast, but it sets the scale at which repeat controls become
necessary for interpretation: a harness-versus-single-call difference of this
magnitude cannot be attributed to capability without them.

\begin{table}[H]
\centering\small
\begin{tabular}{@{}lrr@{}}
\toprule
Sensitivity estimand & Estimate (pp) & 95\% CI (pp) \\
\midrule
K-matched B$-$A headroom & $-0.55$ & $[-1.53, +0.63]$ \\
K-matched D$-$C headroom & $-0.08$ & $[-1.26, +0.80]$ \\
K-matched gate mean & $-0.31$ & $[-1.07, +0.40]$ \\
\midrule
R2 subset D$-$A (cached) & $+0.42$ & $[-0.61, +1.48]$ \\
R2 subset D$-$A (cache-off) & $+0.81$ & $[-0.35, +1.84]$ \\
R2 change in subset D$-$A & $+0.39$ & $[-0.95, +1.63]$ \\
\bottomrule
\end{tabular}

\caption{Corrected sensitivities on core-400, 18 cells. CIs are unadjusted.
K-matched comparisons always retain single-call; R2 comparisons exclude single-call and
are conditional on the observed acquisition groups.}
\label{tab:sensitivity}
\end{table}

\paragraph{Recorded cost and missing billing data.}
The same canonical core records contain \RevisionLogicalCalls{} logged
logical solver calls across \RevisionCostRows{} candidate--task records
(Table~\ref{tab:cost}). These calls include cache-served requests and are
not provider-billed API requests. None of the inspected records has token
usage fields; we do not impute token or monetary expenditure from call
counts. Actual generation attempts, admitted candidate counts, and these
execution counts therefore remain separate measurements. The population
cost mean gives each builder--seed cell equal weight and excludes single-call;
zero-candidate populations would contribute zero solver calls.

\begin{table}[H]
\centering\small
\begin{tabular}{@{}lrrr@{}}
\toprule
Arm & Candidate-task records & Logged solver calls & Calls/task/population \\
\midrule
A & 41,600 & 67,258 & 9.34 \\
B & 44,000 & 64,399 & 8.94 \\
C & 40,400 & 75,604 & 10.50 \\
D & 27,600 & 52,043 & 7.23 \\
\bottomrule
\end{tabular}

\caption{Descriptive canonical cost audit on core-400, 18 cells per arm.
Logged logical calls are not billable requests or token counts. The last
column averages candidate-population calls per task over cells.}
\label{tab:cost}
\end{table}


\FloatBarrier
\subsection{Independent controls for the admission instrument}\label{app:gatecontrols}
We tested the post-review Gate-v2 development instrument with 32 newly
authored controls, separate from the controls used to develop its probes.
The instrument was fixed before a separate model
(\texttt{gpt-5.5}; Table~\ref{tab:models}) received the public
contracts and callable interface, without its implementation or earlier
controls. The author supplied two intended conforming and two intended
nonconforming implementations for each of the eight mechanism types.
Control source, author-reference labels, instrument, measurement code,
and execution order were fixed before the first evaluation. No benchmark
or provider generation was run, and this first measurement is retained unchanged.

\begin{table}[H]
\centering\small
\begin{tabular}{lrrr}\toprule
Declared mechanism & Structural pass & Structural fail & Review required \\\midrule
Repair & 2 & 2 & 0 \\
Execution-result vote & 2 & 2 & 0 \\
Schema linking & 0 & 4 & 0 \\
Hint guard & 0 & 0 & 4 \\
Two views & 0 & 3 & 1 \\
Decomposition & 1 & 3 & 0 \\
Error classification & 0 & 4 & 0 \\
Format guard & 0 & 0 & 4 \\\midrule
Total & 5 & 18 & 9 \\\bottomrule
\end{tabular}
\caption{First evaluation of 32 independently authored controls by the
fixed development instrument. No case had an execution failure. These
are instrument states, not correctness labels or population error rates;
the author's intended positive/negative labels require separate scrutiny.}
\label{tab:gatecontrols}
\end{table}

The nine review-required cases were retained. A third model
(\texttt{gpt-5.6-sol}) received the
full control source and traces without author-reference labels or
automatic verdicts, for source-contract review of all 32 controls and
semantic adjudication of the four designated mechanism types.
Independent source review supported conformance for one decomposition
control and two error-classification controls that the instrument rejected.
The decomposition stub returned the payload of the first plan step whose
text occurred in a request; retaining the earlier step in the context of
the next request therefore retrieved the earlier payload again. The
classification check removed the exact error text but still counted
an ordinary \texttt{Schema:} heading as a category, making the extracted
category set disagree with the expected error class. These are instrument
failures on these controls, not observed failures of a target model.
Three author-intended positive controls were disputed: two schema-linking
controls retained an unrelated table, and one two-view control batched
its requests and mishandled an empty result. Their original labels and
the review disagreements were retained separately. Among the 29 controls
whose author labels and source review agreed, the instrument rejected
3 of 13 conforming controls; all 16 nonconforming controls were rejected
after the designated semantic review. The instrument therefore did not
pass this finite control check. These counts do not estimate error rates
in a generated population.
Semantic adjudication of the designated 16 controls yielded 7 conforming,
8 nonconforming, and 1 indeterminate case. The indeterminate case already
had a structural failure, which was retained without resolving its semantics.
This small, deliberately balanced control set is not a random sample of
generated harnesses. Passing it would not validate the production
interface or replace the common-raw-pool comparison of admission policies.

\paragraph{Revised instrument and a new independent control set.}
Gate-v3 corrected decomposition request attribution and treated category
keywords and shared request templates as review clues rather than automatic
rejections. We fixed this version before fresh model contexts
(\texttt{gpt-5.5} authoring, \texttt{gpt-5.6-sol} review) authored and
reviewed a new set of 32 controls, again two intended positive and two intended
negative controls per mechanism. The review packet withheld author labels,
automatic verdicts, and instrument-derived attribution metadata. All 32
controls received source review; the 20 controls in the four prompt/formulation
groups plus decomposition also received semantic or data-flow review.
Structural outcomes were 8 passes, 10 failures, and 14 review-required cases,
with no execution failures. Following one recorded evidence clarification,
source review classified 15 controls as conforming and 17 as nonconforming;
semantic review classified 8 as conforming, 11 as nonconforming, and 1 as
indeterminate.

The indeterminate classification control requested an intermediate error
category, but the synthetic interface returned SQL text. Its fallback then
selected a semantics repair action. This establishes a response-type mismatch
in the fixture, not a real solver's classification error. A second classification
control was reference-disputed: its source mapped the supplied schema-failure
message to a semantics action. Among the 31 undisputed references, 14 of 15
conforming controls were accepted and one remained unresolved; all 16
nonconforming controls were rejected. The prespecified finite check therefore
remained unmet because of the unresolved case and insufficient agreed positive
coverage for error classification. One voting control's conformance applied
only to result arbitration; it did not establish task-conditioned candidate
generation. These finite mechanism checks do not validate complete harness
operation or replace the planned common-pool policy comparison.

\paragraph{Runtime response contracts.}
We subsequently tested seven archived ungated, strategy-forced candidates,
selecting the first available candidate per strategy in a fixed ordering,
rather than by performance;
no format-guard candidate was available (Appendix~\ref{app:extraction}).
The 27 synthetic scenarios exposed both candidate and instrument limitations.
The two-view candidate skipped its second execution when the first result was
nonempty, contrary to its original ``execute both'' declaration. In contrast,
the schema-linking candidate requested JSON but received plain text, and
the full plan in decomposition requests prevented the frozen probe from
uniquely attributing the current target.
We retained those observations and tested source-assigned response profiles
in 11 further development scenarios. Dedicated subquestion fields identified
the current target, executable SQL payloads reached the final assembly, and
four errors produced by an in-memory SQLite database elicited the corresponding
syntax or schema repair instructions. Schema-linking JSON changed the explicit
table/column restrictions in the SQL-writing request, although the full schema
remained present. Retaining schema context therefore cannot by itself establish
that the linked subset was ignored. These exposed-candidate checks distinguish
response-interface failures from observable candidate deviations; they do not
replace independent calibration or provide revised benchmark outcomes.

\paragraph{Unified interface and premeasurement reference review.}
Gate-v4 combined the source-assigned response interfaces with the five unchanged
Gate-v3 mechanism probes. A separately authored set of 32 controls
(\texttt{gpt-5.6-sol}) underwent
blinded source review before execution. The author was familiar with earlier
control failures; schema-response and declared voting-parameter mismatches found
in this source review were corrected, with the initial versions retained.
The revised references had two agreed conforming and two agreed nonconforming
controls per mechanism. After freezing their source, profiles, and dependencies,
one measurement covered 124 scenarios without execution exceptions. Structural
results were 4 passes, 6 failures, and 22 review-required cases.

A new review context (\texttt{gpt-5.5}) received all source and actual
traces without reference
labels, source-review decisions, or automatic verdicts. Its first judgments
were 18 conforming and 14 nonconforming. Retaining structural failures in the
prespecified hybrid rule accepted all 16 conforming references but also one
nonconforming hint control, so the finite zero-error criterion was not met.
That control correctly restated existing hints but invented a mandatory sort
order when no hint was supplied. A recorded, targeted evidence clarification
corrected this oversight and a two-view judgment that had treated two samples
from one request as independent formulations. The latter control was already
rejected structurally. The clarified decisions agreed with all 32 references,
but this follow-up does not replace the first blinded result or constitute
a new independent calibration pass. No instrument rerun, provider generation,
or benchmark evaluation was performed in this check.

\paragraph{Scope of the declaration vocabulary.}
These controls cover the eight strategy-forced contracts, not the free arm's
four declarations (repair, vote, twostage, and plain). In the original free
prompt, vote specifies multiple samples and selection but not three-way
majority voting; twostage specifies artifact consumption without a particular
response interface. The plain declaration also specifies no control-flow
mechanism, whereas the free prompt requests a control-flow change. Applying
the eight-strategy gate to these declarations
therefore requires separately fixed semantic rules and an indeterminate policy,
rather than silently treating vote as vote3 or choosing a compatible profile
after observing its result. This source-level scope gap does not estimate
false-rejection rates for generated candidates.


Upstream extraction issues and incomplete provider acquisition records
are retained in Appendix~\ref{app:history}.
